\clearpage\SourceLeaf{451}%
\fancyhead[L]{\scriptsize REMARKS.}
\pdfbookmark[1]{REMARKS.}{bemerkungen-coordinaten}
\begin{center}\textbf{REMARKS.}\end{center}

\indent The dimensions of the Earth found by \textsc{Walbeck} were presumably intended to serve as the basis both for the unit of the coordinates and for the various reductions of the survey.

\textsc{Walbeck} and \textsc{Brummer}. De forma et magnitudine telluris. Aboae 1819 p. 16: ‘The mean degree, or \(\frac{1}{90}\) of the meridian quadrant \(=57009^{t}{,}76\). Ellipticity \(=\frac{1}{302{,}78}\)’ [Handwritten remark by \textsc{Gauss}: mean degree of the meridian] \(=\) ‘\(57009^{t}{,}7584\). Hence the metre \(=443^{l}{,}307885\), ratio \(=37299:37300\) logarithm \(=0{,}00001164\).’

\textsc{Gauss}. Bestimmung des Breitenunterschiedes zwischen den Sternwarten von Göttingen und Altona. Göttingen 1828. Art. 20. — ‘If one calculates my triangles on the assumption that they lie on the surface of an elliptical spheroid with the dimensions that \textsc{Walbeck} derived from all the arc measurements made up to that time, a spheroid which, to the best of our present knowledge, fits the actual shape \textit{as a whole} most closely (flattening \(\frac{1}{302{,}76}\), one three-hundred-and-sixtieth of the Earth’s meridian \(=57009{,}746\) toises), and takes the altitude of the pole at Göttingen \(=51^{0}\ 31'\ 37''85\) as the starting point’..

It therefore appears that \textsc{Gauss} repeatedly calculated with the flattening \(\frac{1}{302{,}76}\) in place of \textsc{Walbeck}’s \(\frac{1}{302{,}78}\). Indeed, the former number also underlies several of the auxiliary tables still in the manuscript estate.

\textsc{Gauss} to \textsc{Schumacher}. Göttingen 1830 April 18 ‘Second auxiliary table, note: ‘In the coordinates I communicated earlier, the unit is \(\frac{1}{10000000}\) of the Earth’s meridian quadrant according to \textsc{Walbeck}’s dimensions. To convert them to coordinates based on \(\frac{1}{10000000}\) of the Earth’s meridian quadrant according to \textsc{Schmidt}’s latest, they must first be multiplied by \(\frac{57009758}{57008551}\) or by \(1+\frac{1}{47245}\).’

\textsc{Gauss}. Bestimmung des Breitenunterschiedes zwischen Göttingen und Altona. Art. 19. ‘The trigonometric connection between the Göttingen and Altona observatories places Altona 115163,725 toises north and 7,211 toises west of Göttingen. These numbers refer to the sites of the meridian circles. They rest on the value of 13841,815 toises for the Hamburg-Hohenhorn triangle side, which in turn rests on the baseline that Prof. \textsc{Schumacher} measured in Holstein in 1820. The comparison of the measuring rods used for this purpose with the standard toise has yet to be completed \textit{definitively}. The distance above will therefore have to be adjusted in the same proportion as the baseline itself, although the change can in any event be only very small.’

Privy Councillor of State \textsc{Andrae} in Copenhagen comments on the revision of the baseline in a letter dated 5 March 1865, printed in the Generalbericht über die mitteleuropäische Gradmessung für das Jahr 1864, pp. 6. 7. ‘The length of the \textit{Braak baseline} given by \textsc{Schumacher}: 3014,5799 toises, which was used in the earlier computations of both the Danish triangulation and the Hanoverian triangulation carried out under \textsc{Gauss}’s direction, could only be regarded as a provisional result of the baseline measurement, because the reduction to sea level and several other corrections had yet to be taken into account. These reductions are considerably larger than the uncertainties of the measurements themselves, which were carried out with great care. A new determination was therefore necessary, and Professor Dr. \textsc{Peters} in Altona was kind enough to make a detailed calculation of all the corrections with the greatest accuracy. The length of the baseline is now as follows:
